% Solutions/hints for ch04 exercises (assembled into Appendix C)

\begin{solution}{exr:ch04-trace}
Step~11 expands $(3,4)$ with $\gcost=7$, $\hcost=4$, $\fcost=11$, and \Open becomes $(0,4){:}11$, $(1,5){:}13$, $(2,5){:}13$, $(3,3){:}11$, $(3,5){:}13$. Step~12 expands $(3,3)$ with $\gcost=8$, $\hcost=3$, $\fcost=11$, and \Open becomes $(0,4){:}11$, $(1,5){:}13$, $(2,5){:}13$, $(3,2){:}11$, $(3,5){:}13$. The cell $(0,4)$ has $\fcost=11$ as well, so it ties with the cells of the column $x=3$ at every one of these steps; the tie-breaking rule of \cref{alg:ch04-astar} prefers the larger $\gcost$, and $\gcost(0,4)=4$ is much smaller than $\gcost(3,4)=7$, $\gcost(3,3)=8$, $\gcost(3,2)=9$. The search therefore finishes the dive into the dead end first and pops $(0,4)$ at step~14, immediately after $(3,2)$, whose only successor $(3,1)$ has $\fcost=13$.
\end{solution}

\begin{solution}{exr:ch04-fourheuristics}
At $(1,4)$: $d_x=4$, $d_y=2$, so $\hcost_{\mathrm{M}}=6$, $\hcost_{\mathrm{O}}=4+2(\sqrt2-1)\approx4.83$, $\hcost_{\mathrm{E}}=\sqrt{20}\approx4.47$, $\hcost_{\mathrm{C}}=4$. At $(3,1)$: $d_x=2$, $d_y=1$, so $\hcost_{\mathrm{M}}=3$, $\hcost_{\mathrm{O}}=2+(\sqrt2-1)\approx2.41$, $\hcost_{\mathrm{E}}=\sqrt5\approx2.24$, $\hcost_{\mathrm{C}}=2$. (a) The order is the same everywhere, $\hcost_{\mathrm{M}}\ge\hcost_{\mathrm{O}}\ge\hcost_{\mathrm{E}}\ge\hcost_{\mathrm{C}}$, so the Manhattan distance is the most informed of the four; the first inequality holds because $d_x+d_y=\max+\min\ge\max+(\sqrt2-1)\min$. (b) On the 4-connected instance all four are admissible (\cref{tab:ch04-heuristics}), and the Manhattan distance is exact on an empty grid. On an 8-connected instance the Manhattan distance is not admissible, the other three are. (c) The best of the four is the Manhattan distance, and it underestimates by $8-6=2$ at $(1,4)$ and by $5-3=2$ at $(3,1)$. Both errors come from the second wall, the cells $(4,1)$ to $(4,4)$: from $(1,4)$ the route must climb to the top row before it can cross $x=4$, and from $(3,1)$ it must drop to $y=0$, in each case two moves that the straight-line count does not see.
\end{solution}

\begin{solution}{exr:ch04-octile}
Write $M=\max(d_x,d_y)$ and $m=\min(d_x,d_y)$, so $\hcost_{\mathrm{O}}=M+(\sqrt2-1)m$. On an obstacle-free 8-connected grid a route from $n$ to $\gamma$ needs at least $M$ moves, and each move changes $\min$ by at most one, so at most $m$ of them can be diagonal; taking $m$ diagonal and $M-m$ straight moves is feasible and costs $\sqrt2\,m+(M-m)=M+(\sqrt2-1)m$, which is therefore exactly $\hcost^*$ on an empty grid. For consistency, take a move from $n$ to $n'$ and let $M',m'$ be the values at $n'$. A straight move changes exactly one of $d_x,d_y$ by $\pm1$; a diagonal move changes both by $\pm1$. Enumerating: a straight move that decreases $M$ gives $\hcost_{\mathrm{O}}(n)-\hcost_{\mathrm{O}}(n')=1\le 1$; a straight move that decreases $m$ gives $\sqrt2-1\le1$; a diagonal move that decreases both gives $1+(\sqrt2-1)=\sqrt2\le\sqrt2$; a diagonal move that decreases one and increases the other gives $|1-(\sqrt2-1)|=2-\sqrt2\le\sqrt2$. Moves that increase $M$ or $m$ only make the left-hand side smaller. In every case $\hcost_{\mathrm{O}}(n)\le c(n,n')+\hcost_{\mathrm{O}}(n')$, and $\hcost_{\mathrm{O}}(\gamma)=0$. Finally, consistency is a condition on the edges that exist: forbidding the corner-cutting diagonals of \cref{def:ch02-grid} removes edges, and removing an edge cannot violate an inequality that held on all of them. (Admissibility with obstacles follows because obstacles only lengthen paths.) The same argument is \cref{thm:ch04-metric} with $d$ the octile metric.
\end{solution}

\begin{solution}{exr:ch04-wbound}
Let $\hcost$ be consistent and let no closed node ever be re-opened. Claim: every node $n$ that is expanded satisfies $\gcost(n)\le w\,\gcost^*(n)$. Induct on the order of expansion. For $s$ it is clear. Let $n$ be expanded and let $P=(s,n_1,\dots,n_k=n)$ be a cheapest path to $n$; let $n_i$ be the last node of $P$ that was expanded before $n$ with $\gcost(n_i)\le w\,\gcost^*(n_i)$ (it exists: $s$ qualifies). Then $n_{i+1}$ is in \Open with $\gcost(n_{i+1})\le\gcost(n_i)+c(n_i,n_{i+1})\le w\,\gcost^*(n_{i+1})$. Since $n$ was popped before $n_{i+1}$,
\[\gcost(n)+w\,\hcost(n)\le\gcost(n_{i+1})+w\,\hcost(n_{i+1})\le w\,\gcost^*(n_{i+1})+w\,\hcost(n_{i+1})\le w\bigl(\gcost^*(n_{i+1})+\hcost^*(n_{i+1})\bigr),\]
using admissibility in the last step, and consistency along the rest of $P$ gives $\hcost(n_{i+1})\le\gcost^*(n)-\gcost^*(n_{i+1})+\hcost(n)$, whence $\gcost(n)\le w\,\gcost^*(n)$. Applying the claim to the goal gives $\gcost(\gamma)\le w\,C^*$. Consistency is used exactly once, in the step that propagates the bound along $P$; with a merely admissible $\hcost$ that step fails, a node may be expanded with a $\gcost$ worse than $w\,\gcost^*$, and, with re-opening disabled, that error is never corrected. This is the argument of Likhachev, Gordon and Thrun for \arastar (\cref{ch:ch06}).
\end{solution}

\begin{solution}{exr:ch04-max}
(a) $\hcost(n)=\max(\hcost_1(n),\hcost_2(n))\le\hcost^*(n)$ because each term is, so $\hcost$ is admissible. If both are consistent, then for an edge $(n,n')$ and for $j$ the index attaining the maximum at $n$, $\hcost(n)=\hcost_j(n)\le c(n,n')+\hcost_j(n')\le c(n,n')+\hcost(n')$; and $\hcost(\gamma)=0$. (b) $\hcost\ge\hcost_1$ and $\hcost\ge\hcost_2$ by definition, so \cref{thm:ch04-dominance} says that \astar with $\hcost$ expands, among nodes with $\fcost<C^*$, a subset of what either of the two expands. (c) On the grid of \cref{ex:ch04-grid}, take $\hcost_1$ the Manhattan distance and $\hcost_2(n)=\abs{x_n-5}+2$ on the cells with $x_n\le3$ and $y_n\ge1$ (which must first reach the row $y=0$ or the row $y=5$ to cross the wall at $x=4$, so two extra vertical moves are unavoidable), with $\hcost_2=\hcost_1$ on all other cells: their maximum is strictly larger than $\hcost_1$ at, for example, $(3,1)$, where it gives $4$ instead of $3$. One does not simply maximise everything because each term must be computed at every generated node: a heuristic that is slightly more informed but ten times more expensive loses, and this trade-off between node quality and node rate is the reason the cheap closed-form distances of \cref{tab:ch04-heuristics} remain the default.
\end{solution}

\begin{solution}{exr:ch04-lattice}
Sort the absolute differences as $d_{(1)}\ge d_{(2)}\ge d_{(3)}$. \emph{6-connected} (only axis moves, cost 1): $\hcost_6=d_x+d_y+d_z$, the 3D Manhattan distance, exact on an empty lattice. \emph{26-connected} (costs $1,\sqrt2,\sqrt3$): use $d_{(3)}$ triple moves, then $d_{(2)}-d_{(3)}$ double moves, then $d_{(1)}-d_{(2)}$ single moves, so $\hcost_{26}=\sqrt3\,d_{(3)}+\sqrt2\,(d_{(2)}-d_{(3)})+(d_{(1)}-d_{(2)})$; no route can do better, because each move reduces the sum $d_x+d_y+d_z$ by at most the number of coordinates it changes and $\sqrt3<\sqrt2+1<3$. \emph{18-connected} (axis and face-diagonal moves, costs $1,\sqrt2$): with $S=d_x+d_y+d_z$ the number of double moves is at most $m=\min\bigl(d_{(2)}+d_{(3)},\lfloor S/2\rfloor\bigr)$, and using $m$ of them is feasible, so $\hcost_{18}=\sqrt2\,m+(S-2m)$. Since a richer move set can only make paths cheaper, the true costs satisfy $C_6\ge C_{18}\ge C_{26}$ and the heuristics satisfy $\hcost_6\ge\hcost_{18}\ge\hcost_{26}$ with equality to the corresponding true cost on an empty lattice. Hence $\hcost_{26}$ is admissible on all three lattices, $\hcost_{18}$ on the 6- and 18-connected ones but not on the 26-connected one, and $\hcost_6$ only on the 6-connected one --- the three-dimensional version of the Manhattan trap of \cref{tab:ch04-heuristics}. The Euclidean distance is admissible on all three and is dominated by all of them.
\end{solution}

\begin{solution}{exr:ch04-manhattan8}
Take the $4\times4$ grid whose only obstacles are $(0,1)$ and $(2,1)$, with $s=(0,0)$ and $\gamma=(3,3)$, 8-connected with the corner-cutting rule of \cref{def:ch02-grid}. \astar with the Manhattan heuristic returns $(0,0),(1,0),(2,0),(3,0),(3,1),(3,2),(3,3)$, six straight moves of cost~6, while the optimum is $4+\sqrt2\approx5.41$; \code{dijkstra\_grid(grid, s, 8)} and \astar with the octile heuristic both return $5.41$. The proof of \cref{thm:ch04-optimal} fails at the inequality $\fcost(n_i)=\gcost^*(n_i)+\hcost(n_i)\le\gcost^*(n_i)+\hcost^*(n_i)\le C^*$: with $\hcost>\hcost^*$ the node $n_i$ of the optimal path can have $\fcost(n_i)>C^*$, so nothing forces the algorithm to pop it before the suboptimal goal entry, and the contradiction disappears.
\end{solution}

\begin{solution}{exr:ch04-coding}
Checkpoints rather than a solution. (a) If your expansion count is 22 or 23 instead of 21, you are probably testing the goal at generation time or breaking ties towards the smaller $\gcost$ (24 expansions). If \python raises \code{TypeError: '<' not supported between instances of ...}, you forgot the counter in the priority tuple. (b) The frontier picture is easiest as a grid of coloured squares redrawn after every pop; keep the closed cells grey, the open cells blue, and the current node highlighted, exactly as in \cref{fig:ch04-example}. (c) The three traps are the goal test (\cref{def:ch04-goalstay}), the edge constraint (it is indexed by the departure time $t$, not the arrival time $t+1$) and the horizon: without it your search never terminates on an unsolvable instance. (d) The four assertions of the chapter's self-test are the ones to copy: equal cost to Dijkstra, zero re-openings with a consistent heuristic, a non-decreasing sequence of $\fcost$ values, and \code{space\_time\_astar} cost equal to \code{static\_distances} when there are no constraints. Compare your numbers with \code{ch04\_astar.py} only after your own version runs.
\end{solution}

\begin{solution}{exr:ch04-spacetime}
(a) With both $\langle a,(1,1),1\rangle$ and $\langle a,(1,1),2\rangle$ the middle cell is unusable at $t=1$ and $t=2$, so the agent waits twice: the path is $(0,1),(0,1),(0,1),(1,1),(2,1)$, cost~4, two waits, and \code{space\_time\_astar} expands 5 states. A detour through $(1,0)$ or $(1,2)$ also costs~4 but arrives no earlier, and the tie-breaking rule prefers the state with the larger $\gcost$, hence the wait. (b) The single edge constraint $\langle a,(1,1),(2,1),1\rangle$ forbids the last move at the time the unconstrained plan makes it, so the agent moves to $(1,1)$ at $t=1$, waits there and steps to the goal at $t=3$: the path is $(0,1),(1,1),(1,1),(2,1)$, cost~3, 4~expansions. With the time index~2 instead of~1 the constraint forbids a move the agent never makes at that time, and the unconstrained path $(0,1),(1,1),(2,1)$ of cost~2 stays optimal --- the off-by-one that hides in every hand-written constraint. (c) The expansion counts are 5, 4 and 3 against the 4 of \cref{tab:ch04-spacetime}. (d) With \code{max\_time = 3} instance~(a) is reported unsolvable, because its solution arrives at $t=4$; instance~(b) still returns its 3-step path. This does not contradict \cref{thm:ch04-horizon}, which guarantees completeness only for $T_{\max}=H+1+D$: here $H=2$ and $D=3$, so the sufficient horizon is 6, and 3 is a hand-supplied truncation.
\end{solution}

\begin{solution}{exr:ch04-horizon}
(a) \Cref{ex:ch04-spacetime} itself: $H=1$, so a horizon of $H+1=2$ makes the search fail, while the optimal constrained path arrives at $t=3$ (this is the assertion \code{max\_time=2} $\to$ failure, \code{max\_time=3} $\to$ cost 3 in the self-test). (b) The same instance in the other direction: $T_{\max}=1+1+3=5$ while the arrival time is 3. Any instance whose constraints are far from the route makes the gap arbitrarily large, because $D$ is the eccentricity of the goal and does not look at the start. (c) For agent~2 the table contains agent~1's five vertex constraints, its four reversed edges and the parked cell $(4,0)$ from $t=4$, so $H=4$; treating $(4,0)$ as a wall, the largest distance to $(0,0)$ is $D=3$ (from the bay $(2,1)$), giving $T_{\max}=8$, while the search arrives at $t=5$. (d) Take the ring map: the $4\times4$ grid whose four interior cells are blocked, so the 12 border cells form a cycle. Let $\gamma=(0,0)$, let the agent start at $(2,0)$ and let another agent be parked on $(1,0)$ from $t=0$, so $H=0$. In the raw grid the eccentricity of $\gamma$ is 6, giving the naive horizon $0+1+6=7$; but the parked cell cuts the short side of the ring and the only route left is the long way round, ten moves, so the true arrival time is 10. With \code{max\_time = 7} the search returns failure, and a prioritized planner (\cref{ch:ch08}) would report an unsolvable instance that in fact has a solution. Computing $D$ in the grid \emph{with} the parked cells treated as walls, as \code{default\_horizon()} does, gives $D=10$ and $T_{\max}=11$, and the path is found.
\end{solution}
